% !TeX root = ../../../main.tex
\section{Introduction}

One in 10 people globally are affected by chronic kidney disease (CKD) \cite{bikbov_Global_2020}, a comorbidity to a wide range of chronic etiologies including inherited diseases such as polycystic kidney disease, acquired conditions including diabetes and hypertension, and recurrent renal infections. The end-stage variant of CKD can be addressed in the most efficacious and cost-effective way through renal allografts. While only 2.9\% of renal transplants performed in the United States of America between 2010 and 2013 have failed in less than one year, 32.4\% of them failed within a decade \cite{poggio_Longterm_2021}. When renal graft dysfunction is suspected, no definitive treatment currently exists for AMR~\cite{bohmig_Antibodymediated_2025}, and the long-term efficacy of existing TCMR treatments remain unclear~\cite{ho_Effectiveness_2022, rampersad_negative_2022}. Furthermore, attention has recently been drawn to the importance of "subclinical rejection"~\cite{loupy_Subclinical_2015,mehta_Longterm_2022}, the presence of histological findings of rejection in the absence of significant renal dysfunction. Together, this suggests that the proactive monitoring and early detection of renal allograft pathologies is more critical than ever for improving long-term graft survival.

Failures of renal allografts may be broadly attributed to two categories of causes: those that occur due to immunologic causes and those due to non-immunologic causes~\cite{saritas_Kidney_2021}. Immunologic causes include T-cell mediated rejection (TCMR) and antibody-mediated rejection (AMR), where the recipient's immune system mounts a response against the grafted tissue. Non-immunologic causes include acute kidney injury, cytotoxicity from immunosuppressants (e.g. calcineurin inhibitors, or CNI), infections, and recurrent primary disease~\cite{saritas_Kidney_2021,naesens_Banff_2024}. While the specific mechanisms of these causes differ, many of them manifest as abnormalities that can be observed in needle core biopsies and characterized using the Banff Classification of Renal Allograft Pathology~\cite{loupy_Banff_2020,naesens_Banff_2024}. The Banff system defines specific pathologies using a combination of proteomic evidence such as the C4d complement fragment in the renal parenchyma, serological abnormalities such as the circulation of donor-specific antibodies (DSA) for human leukocyte antigen (HLA), and findings of histological changes in the renal parenchyma~\cite{alexandre_AntibodyMediated_2018}. Specific lesion types are graded on ordinal scales by severity and contextualized to diagnose TCMR, AMR, interstitial fibrosis and tubular atrophy (IFTA), and other specific disease entities~\cite{naesens_Banff_2024}. Notably, the final diagnoses are not mutually exclusive. IFTA~\cite{saritas_Kidney_2021}, as well as observations of interstitial inflammation~\cite{mannon_Inflammation_2010}, have been associated with lower graft survival rates. Furthermore, the co-localized presence of inflammation in areas of IFTA has been shown to be particularly prognostically significant~\cite{gago_Kidney_2012,nankivell_causes_2018}.

The current gold standard for diagnosing renal allograft pathologies remains the histological analysis of needle core biopsies~\cite{naesens_Banff_2024}. However, these biopsies are invasive procedures that carry non-negligible risks of complications including hematuria, hematoma formation, and arteriovenous fistula formation to the point where transfusions or surgical interventions may be warranted~\cite{geddes_Major_2025,ho_Complications_2022,morgan_Complications_2016}. Therefore, the regular use of the aforementioned needle biopsies for renal allograft monitoring is controversial among nephrologists~\cite{schinstock_Value_2017}. While some institutions and practitioners value the wealth of cell-level information that biopsies provide, others restrict their usage to patients showing other pathological indications to avoid patient discomfort and undue risks for injury or infection.

However, alternative, noninvasive methods for monitoring renal allograft health remain limited. Serum biomarkers such as serum creatinine and blood urea nitrogen (BUN), as well as urine-based biomarkers such as proteinuria, estimated glomerular filtration rate (eGFR), and urine albumin-to-creatinine ratio (UACR), are widely used due to their low cost and ease of measurement. However, all of those approaches lack sensitivity for detecting renal pathologies, and are nonspecific to the etiology of graft dysfunction~\cite{goutaudier_Evaluation_2024,levey_Strengths_2019,naesens_Proteinuria_2016,porrini_Estimated_2019,sasaki_Variability_2025}. DSA monitoring adds specificity for detecting AMR~\cite{haas_Differences_2017}, but its adoption is variable~\cite{vandenbroek_Clinical_2023} and it does not address TCMR, non-immunologic causes of graft dysfunction, and AMR in the absence of circulating DSAs~\cite{loupy_Subclinical_2015}. Conventional imaging (B-mode ultrasound, MRI, scintigraphy) provides anatomical and some functional information~\cite{sugi_Imaging_2019,berchtold_Diffusionmagnetic_2022}, yet it is constrained by operator dependence, hardware access, and limited ability to encode microstructural changes that mirror Banff lesion scoring.

% Second half of above paragraph used to be its own separate paragraph:
%
% Imaging modalities such as conventional B-mode ultrasound and magnetic resonance imaging (MRI) can provide anatomical information regarding renal allograft morphology and gross abnormalities such as hydronephrosis or vascular occlusions~\cite{sugi_Imaging_2019}. Beyond its noninvasiveness, imaging-based approaches are also useful for its multimodality; functional information regarding bloodflow and filtration may be added to improve the ability to detect renal graft dysfunction. Additional information may be gained through Doppler imaging or the use of contrast agents for ultrasound, while diffusion-weighted imaging (DWI) and blood oxygenation labeling (BOLD) can be used using MRI~\cite{berchtold_Diffusionmagnetic_2022,sugi_Imaging_2019,irshad_review_2008}. However, the ultrasound-based techniques are limited by operator skills and poor repeatability, while MRI- and scintigraphy-based techniques are limited by access to hardware, long acquisition times, and contraindications in certain patient populations. Plus, both functional ultrasound and MRI techniques do not encode microstructural information regarding the renal parenchyma such that they may serve as an adjunct or alternative to biopsies for detecting graft pathologies.

Elastography, or the imaging of tissue mechanical properties, has recently emerged as a promising noninvasive tool for assessing renal allograft health. Several techniques that involve the emission of acoustic radiation force (ARF), or the transfer of momentum from an ultrasound wave to tissue, have been proposed for renal allograft assessment~\cite{urban_Novel_2021}. Many of these approaches have involved the tracking of shear waves induced by ARF excitations, where the shear wave speed is related to the shear modulus of the tissue~\cite{urban_Novel_2021,filipov_Investigating_2024}. However, shear wave-based techniques such as shear wave elasticity imaging (SWEI or SWE) and shearwave dispersion ultrasound vibrometry (SDUV) require the assumption that tissue is homogeneous and, unless carefully controlled for, isotropic (i.e. not dependent on the direction of measurement) throughout the entire region where shear waves propagate~\cite{sarvazyan_Shear_1998,ngo_plane_2024,andrade_Simultaneous_2025}. These assumptions are often violated in renal tissue due to its complex microstructure, which includes anisotropic arrangements of tubules and vessels that affect shear wave propagation. Furthermore, shear wave-based techniques are sensitive to reflections and interferences from anatomical boundaries, which are abundant in the kidney due to its heterogeneous structure~\cite{urban_Novel_2021}. These limitations may confound the ability of shear wave-based techniques to accurately assess renal allograft health. As such, alternative ARF-based elastography techniques that do not rely on shear wave propagation may provide complementary information regarding renal allograft health.

Viscoelastic response (VisR) ultrasound, introduced by Selzo and Gallippi~\cite{selzo_Viscoelastic_2013} as an ARF-based method that evaluates elasticity and viscosity at the site of an ARF excitation, may offer such an alternative. Rather than tracking shear waves over a distance and measuring the speed of wave propagation, VisR tracks the displacement of tissue on-axis to an ARF excitation and uses the resulting displacement profile to estimate tissue viscoelastic properties~\cite{hossain_Viscoelastic_2020}. The VisR-derived metrics of regional ratios (RR) and degree of anisotropy (DoA) have been shown to be effective in distinguishing viscoelastic differences in pathological renal allografts using an acquisition protocol that is less sensitive to reflections and interferences from anatomical boundaries~\cite{hossain_Evaluating_2018,hossain_Quantitative_2022}. Previous works have demonstrated that VisR-based DoA in healthy transplant recipients did not vary in a significantly different manner according to recipients' genders, races, donor sources, or BMI~\cite{hossain_Mechanical_2019}. It has also been shown that changes in elastic anisotropy trend with DoA based on ARFI PD \emph{in silico} and through experimental acquisitions~\cite{hossain_Viscoelastic_2020,hossain_Quantitative_2022}. However, the assessment of VisR-based RR and DoA for the early detection of renal allograft pathologies \emph{in vivo}, as well as their direct comparisons against Banff histological findings, have not yet been reported.

As such, we hypothesized that VisR can identify renal allografts with microstructural pathologies\emph{in vivo} through RR and DoA measurements. The performance of VisR-specific imaging metrics were compared against ARFI peak displacement (PD) as well as clinical biomarkers of renal function and histological findings from indication biopsies.
